% from main.tex
%%%%%%%%%%  scan idx 537  |  folio 524  %%%%%%%%%%
\[
x\cl_\bullet(\mathcal O_Z)\in\Filt_{j-i}(X).
\]
Pour éviter toute confusion, désignons par $1_Z$ la classe de $\mathcal O_Z$ dans $K_\bullet(Z)$. Si $\alpha:Z\to X$ est l'immersion canonique, alors la formule de projection (IV \ldots) donne
\[
x\cl_\bullet(\mathcal O_Z)=x\alpha_*(1_Z)=\alpha_*(\alpha^*(x)1_Z).
\]
On a $\alpha^*(x)\in\Filt^i(Z)$, donc d'après (1.1.8) et (1.3.3) pour $X=Z$, $i+j\leq k$, on trouve
\[
\alpha_*(\alpha^*(x)1_Z)\in\Filt_{j-i}(X),
\]
d'où l'assertion.

Le théorème 1.3.2 est trivial pour $i+j\leq0$. Supposons-le démontré par récurrence pour $i+j<k$ et pour tout schéma noethérien $X$. D'après la remarque précédente il nous reste à démontrer que pour tout schéma noethérien intègre $X$ de dimension $j$ et pour tout produit (1.2.3) satisfaisant à (1.2.4), on a
\begin{equation}\tag{1.3.4}
P\cl_\bullet(\mathcal O_X)\in\Filt_{j-i}(X).
\end{equation}
Si $k>1$ dans (1.2.3), on a gagné grâce à l'hypothèse de récurrence et l'associativité de la loi (1.3.1). On peut donc supposer ($X$ étant connexe) que
\begin{equation}\tag{1.3.5}
x=P=\gamma^m(\cl^\circ(E)-n),
\end{equation}
où $m\geq i$ et $E$ est un $\mathcal O_X$-module localement libre de rang $n>0$. Distinguons deux cas :

I. Supposons qu'on peut trouver des $\mathcal O_X$-modules inversibles $L_1,\ldots,L_n$ tels que
\[
\cl^\circ(E)=\cl^\circ(L_1)+\cdots+\cl^\circ(L_n)\in K^\circ(X).
\]
